% appendix_sink.tex -- \input by paper.tex after \appendix

\section{Proofs}
\label{app:proofs}

Throughout, $D$ is a symmetric $N\times N$ matrix with non-negative entries and zero diagonal, $K_\epsilon$ is the Vietoris--Rips complex at scale $\epsilon$, $\mathrm{filt}_D(\sigma)=\max_{u,v\in\sigma}D_{uv}$ (with $\mathrm{filt}_D(\{u\})=0$), and homology is taken over a field. Bars are half-open intervals $[b,d)$.

\subsection{Fact~\ref{fact:mst}}
All vertices enter at $0$. Processing edges in non-decreasing weight order, an edge either joins two components, in which case (elder rule) one $0$D class dies at its weight and Kruskal's algorithm adds the edge to the spanning forest, or it closes a cycle, in which case no $0$D class dies and Kruskal discards it. Hence the $N-1$ finite death times are the MST edge weights \citep{gower1969,edelsbrunner2010}.

\subsection{Theorem~\ref{thm:cone}}

\begin{lemma}[Coned filtrations]
\label{lem:coned}
If $D$ is coned at $s$, then for every $\epsilon\ge0$, $K_\epsilon=C_\epsilon\sqcup I_\epsilon$, where $C_\epsilon$ is a simplicial cone with apex $s$ and $I_\epsilon$ is a set of isolated vertices. Consequently $\tilde H_k(K_\epsilon)=0$ for every $k\ge1$.
\end{lemma}
\begin{proof}
Let $\sigma\in K_\epsilon$ with $s\notin\sigma$ and $|\sigma|\ge2$. For each $u\in\sigma$ choose $v\in\sigma\setminus\{u\}$; coning gives $D_{su}\le D_{uv}\le\mathrm{filt}_D(\sigma)\le\epsilon$. Therefore $\mathrm{filt}_D(\sigma\cup\{s\})=\mathrm{filt}_D(\sigma)$ and $\sigma\cup\{s\}\in K_\epsilon$. Now let $u\ne s$ be a vertex with $D_{su}>\epsilon$. If $u$ had an edge $\{u,v\}\in K_\epsilon$ then $D_{su}\le D_{uv}\le\epsilon$, a contradiction; so $u$ is isolated. Let $I_\epsilon$ be these vertices and $C_\epsilon$ the subcomplex on the remaining vertices. Every vertex of $C_\epsilon$ is adjacent to $s$, and by the first part every simplex of $C_\epsilon$ extends by $s$, so $C_\epsilon$ is the cone over its link of $s$ and is contractible. $H_k$ of a disjoint union is the direct sum, and isolated points have no homology in degree $\ge1$.
\end{proof}

\begin{proof}[Proof of Theorem~\ref{thm:cone}]
Fix $s$ and $\defect=\defect_s(D)$. Define $D'_{uv}=\max(D_{uv},D_{su},D_{sv})$ for $u,v\ne s$, $u\ne v$; $D'_{su}=D'_{us}=D_{su}$; $D'_{uu}=0$.

\emph{Step 1: $D\le D'\le D+\defect$ and $D'$ is coned at $s$.} The first inequality is immediate. For $u,v\ne s$, $D'_{uv}-D_{uv}=[\max(D_{su},D_{sv})-D_{uv}]_+\le\defect$. Coning: $\max(D'_{su},D'_{sv})=\max(D_{su},D_{sv})\le D'_{uv}$.

\emph{Step 2: interleaving.} Since filtration values are maxima of entries, $\mathrm{filt}_D\le\mathrm{filt}_{D'}\le\mathrm{filt}_D+\defect$. Hence for every $b$, $K_b\subseteq K'_{b+\defect}\subseteq K_{b+\defect}$, where $K'$ is the filtration of $D'$.

\emph{Step 3: (a).} The persistence map $H_k(K_b)\to H_k(K_{b+\defect})$ is induced by the composite inclusion in Step 2 and therefore factors through $H_k(K'_{b+\defect})$, which is zero for $k\ge1$ by Lemma~\ref{lem:coned}. The rank of $H_k(K_b)\to H_k(K_{b'})$ equals the number of bars $[b_i,d_i)$ with $b_i\le b$ and $d_i>b'$. A bar with $d_i-b_i>\defect$ would contribute to this rank at $b=b_i$, $b'=b_i+\defect$, a contradiction. So every bar has length at most $\defect$, and summing gives the bound on $P_k$.

\emph{Step 4: (b).} By Fact~\ref{fact:mst}, $P_0(D)=w_D(T_D)$ for an MST $T_D$ of $D$. The star $T_s=\{\{s,u\}\}_{u\ne s}$ is a spanning tree, so $P_0(D)\le w_D(T_s)=\stw_s(D)$. For the lower bound, $T_s$ is an MST of $D'$: every non-star edge $\{u,v\}$ closes the cycle $s,u,v$ in which it is a heaviest edge ($D'_{uv}\ge\max(D'_{su},D'_{sv})$), so by the cycle property it can be excluded. Hence $w_{D'}(T_{D'})=w_{D'}(T_s)=\stw_s(D)$, and
$P_0(D)=w_D(T_D)\ge w_{D'}(T_D)-(N-1)\defect\ge w_{D'}(T_{D'})-(N-1)\defect=\stw_s(D)-(N-1)\defect.$
The maximum finite $0$D death is the bottleneck value $\min_T\max_{e\in T}w(e)$, which is attained by any MST; it is at most $\max_u D_{su}$ (the star is a spanning tree) and is $1$-Lipschitz in the sup-norm, which gives the lower bound because the bottleneck of $D'$ is $\max_u D_{su}$.

\emph{Step 5: (c).} If $\defect=0$ then $D'=D$, and the claims follow from Lemma~\ref{lem:coned} and Step 4.
\end{proof}

\begin{proof}[Proof of Corollary~\ref{cor:causal}]
For $u>v$, $A_{vu}=0$, so $D_{uv}=1-A_{uv}$; in particular $D_{0u}=1-A_{u0}$. For $1\le v<u$, $\max(D_{0u},D_{0v})-D_{uv}=A_{uv}-\min(A_{u0},A_{v0})$. Substitute into Theorem~\ref{thm:cone}.
\end{proof}

\subsection{Proposition~\ref{prop:morse}}
\begin{proof}
With i.i.d.\ uniform weights, the edges with $D_{uv}\le p$ form an Erd\H{o}s--R\'enyi graph $G(N,p)$ and $K_p$ is its clique (flag) complex. Every $1$D bar is contained in $[0,1)$, because $K_1$ is the full simplex. Writing $\beta_1(p)$ for the first Betti number of $K_p$, the number of bars alive at $p$ is $\beta_1(p)$, so $P_1=\int_0^1\beta_1(p)\,dp$.
For any finite complex with $f_i$ faces of dimension $i$, the strong Morse inequality at degree $2$ reads $\beta_2-\beta_1+\beta_0\le f_2-f_1+f_0$ \citep{edelsbrunner2010}, hence $\beta_1\ge f_1-f_0-f_2+\beta_0+\beta_2\ge f_1-f_0-f_2$. Since also $\beta_1\ge0$, Jensen's inequality gives $\mathbb{E}\beta_1(p)\ge[\mathbb{E}f_1-f_0-\mathbb{E}f_2]_+=[\binom N2p-N-\binom N3p^3]_+$. Integrating over $p$ gives the first inequality. With $p=x/\sqrt N$ the integrand is $N^{3/2}(x/2-x^3/6)-N+O(N^{1/2})$, positive on $x\in(2/\sqrt N\cdot(1+o(1)),\sqrt3)$, and
$\int_0^1[\cdot]_+dp=N\int_0^{\sqrt3}\big(\tfrac x2-\tfrac{x^3}{6}\big)dx-O(\sqrt N)=\tfrac38N-O(\sqrt N).$
For the upper bound, $\beta_1\le f_1$ gives $\mathbb{E}P_1\le\int_0^1\binom N2p\,dp=\binom N2/2$.
\end{proof}
Numerically, the integral equals $0.79, 3.92, 11.9, 30.2, 70.1$ at $N=16,32,64,128,256$, against simulated means $1.84, 6.46, 22.4, 64.9, 168.0$ (Figure~\ref{fig:synthetic}, left). Both grow super-linearly over this range (log-log slopes $1.55$ and $1.61$); the lower bound's slope tends to $1$ as $N\to\infty$.

\subsection{Directed flag complexes on causal attention}
\label{app:directed}
\begin{theorem}[Directed flag collapse]
\label{thm:directed}
Let $A$ be lower-triangular with $A_{ij}>0$ for all $j\le i$ (the diagonal included; self-loops are ignored). Let $\mathrm{DFl}(A)$ be the filtered directed flag complex \citep{flagser2020} of the digraph with an edge $i\to j$ for every $j<i$ and weight $1-A_{ij}$, and let $\VR(D)$ be as above with $D=1-\max(A,A^\top)$. Then $\mathrm{DFl}(A)$ and $\VR(D)$ are isomorphic filtered complexes, so $\dgm_k(\mathrm{DFl}(A))=\dgm_k(\VR(D))$ for all $k$. The same holds for any weighting $w(i\to j)=\phi(A_{ij})$ with $D_{ij}=\phi(\max(A_{ij},A_{ji}))$ and $\phi$ non-increasing.
\end{theorem}
\begin{proof}
A directed $k$-simplex is an ordered tuple $(v_0,\dots,v_k)$ with an edge $v_a\to v_b$ for all $a<b$. Edges exist only from later to earlier tokens, so a vertex set $S$ supports exactly one directed simplex, its elements in decreasing order; every other ordering uses a missing forward edge. (A missing edge is not a zero-weight edge; encoding it as weight $0$ would admit all orderings, a common implementation error.) This bijection preserves faces. Its filtration value is $\max_{a<b}(1-A_{v_av_b})$, and for $u>v$, $D_{uv}=1-A_{uv}$, so the undirected value $\max_{u,v\in S}D_{uv}$ is the same number.
\end{proof}
We verified the bijection and the filtration identity by explicit enumeration on $420$ random causal matrices ($N\le8$, four families including tie-dense and underflowing matrices; $30{,}060$ vertex sets; zero violations), and confirmed that the identity fails for bidirectional attention, where a set of size $n$ supports $n!$ orderings. The collapse is a consequence of causality, so directional treatments that are well motivated for bidirectional encoders \citep{kushnareva2021} do not transfer to decoders.

\section{Implementation and protocol details}
\label{app:impl}

\paragraph{Extraction.} Models run in float16 on Apple-silicon (MPS) with eager attention. For each example we run one forward pass with attentions and hidden states, average heads, and form $D$ in float64. Ripser \citep{bauer2021ripser} computes $0$D/$1$D barcodes; per-head MSTs use a batched dense Prim's algorithm (SciPy's sparse MST drops explicit zero-weight edges, which occur when $A_{u0}\approx1$). $\defect_0$ uses the closed form of Corollary~\ref{cor:causal}; $\min_s\defect_s$ is computed by brute force over all apexes. \textsc{Deflated} deletes token 0 and renormalizes each row over the remaining causal support. \textsc{Answer-only} restricts $D$ to the answer tokens without renormalization. The Lookback ratio of head $h$ is the mean over answer tokens $t$ of $\bar a_{\mathrm{ctx}}/(\bar a_{\mathrm{ctx}}+\bar a_{\mathrm{new}})$, with $\bar a$ the mean attention from $t$ to prompt tokens and to answer tokens up to $t$. The LLM-Check score of a head is the mean log diagonal attention (the log-determinant of the triangular attention kernel divided by $N$), summed over heads per layer. Log-probability features are the mean, minimum, and sum of answer-token log-probabilities, mean and maximum predictive entropy, and mean maximum softmax probability. Hidden states are mean-pooled over answer tokens at layer $\lfloor L/2\rfloor$.

\paragraph{Prompts.} TruthfulQA uses each model's native format: the Qwen chat template with the default system turn, \texttt{<|user|>\ldots<|end|>} for Phi-3, \texttt{<|user|>\ldots</s>} for TinyLlama, ChatML without system turn for SmolLM, and \texttt{[INST]\ldots[/INST]} for Mistral. HaluEval uses \texttt{Knowledge: \ldots\textbackslash nQuestion: \ldots\textbackslash nAnswer: }. The on-policy prompt is the model's chat template around ``Answer the following question with a short phrase only, no explanation. Question: \ldots''. Token 0 is whatever token the format places first (BOS for TinyLlama and Mistral, \texttt{<|user|>} for Phi-3, whose tokenizer adds no BOS, \texttt{<|im\_start|>} for Qwen and SmolLM, and ``Knowledge'' for Qwen on HaluEval); the theorem does not require it to be a special token.

\paragraph{On-policy benchmark.}
\begin{table}[h]
\centering
\small
\caption{\textbf{TriviaQA on-policy benchmark.} $2{,}000$ validation questions sampled with seed 0; greedy decoding, first line of the answer, at most 24 new tokens. An answer is correct if any normalized gold alias (lower-cased, punctuation and articles removed) occurs in the normalized answer as a whole-word span. Empty answers are dropped.}
\label{tab:onpolicy}
\SinkTableOnpolicy
\end{table}

\paragraph{Statistics.} Bootstrap resamples are drawn over questions (both rows of a pair move together). TOST power at margin $\epsilon$ assumes a true difference of $0$ and a normal sampling distribution with the bootstrap SE: $2\Phi(\epsilon/\mathrm{SE}-z_{0.95})-1$. Repeated-CV predictions are averaged before bootstrapping, so the intervals reflect sampling of questions for a fixed, split-averaged predictor; the seed-to-seed SD of the pooled AUC is reported separately and is below $0.004$ everywhere.

\section{Full results}
\label{app:tables}
\small
\begin{longtable}{llrrrr}
\toprule
Setting & Bank & dim & AUC & seed SD & pair acc. \\
\midrule\endhead
TQA Qw3B & NONTOPO+DEFL & 2954 & 0.928 & 0.0007 & 0.931 \\
TQA Qw3B & 0D & 72 & 0.754 & 0.0017 & 0.799 \\
TQA Qw3B & SINK & 72 & 0.740 & 0.0028 & 0.765 \\
TQA Qw3B & 1D & 144 & 0.602 & 0.0018 & 0.583 \\
TQA Qw3B & DEFL & 216 & 0.717 & 0.0009 & 0.748 \\
TQA Qw3B & ANS & 216 & 0.752 & 0.0017 & 0.745 \\
TQA Qw3B & ROWSTAT & 72 & 0.751 & 0.0003 & 0.829 \\
TQA Qw3B & LOGPROB & 6 & 0.631 & 0.0008 & 0.682 \\
TQA Qw3B & HIDDEN & 2048 & 0.913 & 0.0008 & 0.916 \\
TQA Qw3B & LEX & -- & 0.820 & 0.0014 & 0.829 \\
TQA Qw3B & LEGACY MSTPROXY & 36 & 0.714 & 0.0006 & 0.786 \\
TQA Qw3B & LLMCHECK & 36 & 0.693 & 0.0011 & 0.764 \\
TQA Qw3B & LOOKBACK & 576 & 0.897 & 0.0026 & 0.906 \\
TQA Qw3B & PH 0D & 1152 & 0.854 & 0.0022 & 0.857 \\
TQA Qw3B & PH 0DTOT & 576 & 0.844 & 0.0010 & 0.868 \\
TQA Qw3B & PH SINK & 576 & 0.850 & 0.0018 & 0.865 \\
TQA Qw3B & PH ENT & 576 & 0.859 & 0.0016 & 0.906 \\
TQA Qw3B & 0D+SINK & 144 & 0.757 & 0.0019 & 0.788 \\
TQA Qw3B & 0D+1D & 216 & 0.760 & 0.0012 & 0.789 \\
TQA Qw3B & SINK+DEFL & 288 & 0.772 & 0.0032 & 0.802 \\
TQA Qw3B & SINK+ANS & 288 & 0.784 & 0.0026 & 0.791 \\
TQA Qw3B & LEX+DEFL & -- & 0.829 & 0.0010 & 0.842 \\
TQA Qw3B & LEX+SINK & -- & 0.843 & 0.0021 & 0.859 \\
TQA Qw3B & LEX+0D & -- & 0.847 & 0.0016 & 0.860 \\
TQA Qw3B & HIDDEN+0D & 2120 & 0.914 & 0.0009 & 0.917 \\
TQA Qw3B & LOGPROB+0D & 78 & 0.787 & 0.0022 & 0.820 \\
TQA Qw3B & SINK+1D & 216 & 0.747 & 0.0021 & 0.767 \\
TQA Qw3B & NONTOPO+0D & 2810 & 0.928 & 0.0004 & 0.936 \\
TQA Qw3B & LEN & 2 & 0.529 & 0.0003 & 0.567 \\
TQA Qw3B & NONTOPO & 2738 & 0.927 & 0.0006 & 0.938 \\
TQA Qw3B & PH SINK+PH 0DTOT & 1152 & 0.860 & 0.0032 & 0.889 \\
TQA Qw1.5B & 0D & 56 & 0.726 & 0.0007 & 0.767 \\
TQA Qw1.5B & LEN & 2 & 0.529 & 0.0003 & 0.567 \\
TQA Qw1.5B & LEX+SINK & -- & 0.833 & 0.0008 & 0.846 \\
TQA Qw1.5B & LEX+0D & -- & 0.835 & 0.0017 & 0.841 \\
TQA Qw1.5B & HIDDEN+0D & 1592 & 0.916 & 0.0014 & 0.924 \\
TQA Qw1.5B & LOGPROB+0D & 62 & 0.740 & 0.0013 & 0.797 \\
TQA Qw1.5B & NONTOPO+DEFL & 2130 & 0.917 & 0.0014 & 0.923 \\
TQA Qw1.5B & NONTOPO+0D & 2018 & 0.917 & 0.0014 & 0.927 \\
TQA Qw1.5B & NONTOPO & 1962 & 0.917 & 0.0013 & 0.928 \\
TQA Qw1.5B & PH SINK+PH 0DTOT & 672 & 0.841 & 0.0013 & 0.867 \\
TQA Qw1.5B & SINK+1D & 168 & 0.715 & 0.0017 & 0.769 \\
TQA Qw1.5B & SINK+ANS & 224 & 0.761 & 0.0027 & 0.753 \\
TQA Qw1.5B & SINK+DEFL & 224 & 0.757 & 0.0028 & 0.802 \\
TQA Qw1.5B & 0D+1D & 168 & 0.718 & 0.0013 & 0.767 \\
TQA Qw1.5B & 0D+SINK & 112 & 0.735 & 0.0012 & 0.766 \\
TQA Qw1.5B & LEX+DEFL & -- & 0.825 & 0.0019 & 0.836 \\
TQA Qw1.5B & PH SINK & 336 & 0.830 & 0.0008 & 0.851 \\
TQA Qw1.5B & PH ENT & 336 & 0.833 & 0.0025 & 0.852 \\
TQA Qw1.5B & 1D & 112 & 0.586 & 0.0014 & 0.644 \\
TQA Qw1.5B & DEFL & 168 & 0.703 & 0.0042 & 0.755 \\
TQA Qw1.5B & ANS & 168 & 0.697 & 0.0040 & 0.700 \\
TQA Qw1.5B & ROWSTAT & 56 & 0.738 & 0.0014 & 0.815 \\
TQA Qw1.5B & LOGPROB & 6 & 0.638 & 0.0007 & 0.682 \\
TQA Qw1.5B & SINK & 56 & 0.723 & 0.0019 & 0.771 \\
TQA Qw1.5B & LEX & -- & 0.820 & 0.0014 & 0.829 \\
TQA Qw1.5B & LEGACY MSTPROXY & 28 & 0.698 & 0.0010 & 0.786 \\
TQA Qw1.5B & LLMCHECK & 28 & 0.666 & 0.0005 & 0.743 \\
TQA Qw1.5B & LOOKBACK & 336 & 0.851 & 0.0020 & 0.856 \\
TQA Qw1.5B & PH 0D & 672 & 0.838 & 0.0012 & 0.847 \\
TQA Qw1.5B & HIDDEN & 1536 & 0.916 & 0.0012 & 0.923 \\
TQA Qw1.5B & PH 0DTOT & 336 & 0.819 & 0.0011 & 0.838 \\
TQA Phi-3 & LOOKBACK & 1024 & 0.892 & 0.0004 & 0.903 \\
TQA Phi-3 & LLMCHECK & 32 & 0.653 & 0.0007 & 0.754 \\
TQA Phi-3 & LEGACY MSTPROXY & 32 & 0.614 & 0.0008 & 0.695 \\
TQA Phi-3 & LEX & -- & 0.820 & 0.0014 & 0.829 \\
TQA Phi-3 & HIDDEN & 3072 & 0.933 & 0.0007 & 0.947 \\
TQA Phi-3 & ANS & 192 & 0.623 & 0.0048 & 0.655 \\
TQA Phi-3 & ROWSTAT & 64 & 0.661 & 0.0011 & 0.734 \\
TQA Phi-3 & DEFL & 192 & 0.639 & 0.0011 & 0.685 \\
TQA Phi-3 & 1D & 128 & 0.529 & 0.0021 & 0.551 \\
TQA Phi-3 & SINK & 64 & 0.652 & 0.0007 & 0.707 \\
TQA Phi-3 & PH 0D & 2048 & 0.868 & 0.0006 & 0.879 \\
TQA Phi-3 & LOGPROB & 6 & 0.578 & 0.0017 & 0.599 \\
TQA Phi-3 & PH 0DTOT & 1024 & 0.846 & 0.0020 & 0.884 \\
TQA Phi-3 & LEX+DEFL & -- & 0.802 & 0.0010 & 0.827 \\
TQA Phi-3 & PH ENT & 1024 & 0.864 & 0.0011 & 0.897 \\
TQA Phi-3 & 0D+SINK & 128 & 0.658 & 0.0007 & 0.707 \\
TQA Phi-3 & 0D+1D & 192 & 0.645 & 0.0012 & 0.707 \\
TQA Phi-3 & SINK+DEFL & 256 & 0.685 & 0.0015 & 0.727 \\
TQA Phi-3 & SINK+ANS & 256 & 0.672 & 0.0013 & 0.709 \\
TQA Phi-3 & SINK+1D & 192 & 0.654 & 0.0010 & 0.705 \\
TQA Phi-3 & NONTOPO & 4198 & 0.936 & 0.0006 & 0.941 \\
TQA Phi-3 & NONTOPO+0D & 4262 & 0.936 & 0.0006 & 0.944 \\
TQA Phi-3 & NONTOPO+DEFL & 4390 & 0.937 & 0.0007 & 0.942 \\
TQA Phi-3 & LOGPROB+0D & 70 & 0.656 & 0.0010 & 0.699 \\
TQA Phi-3 & HIDDEN+0D & 3136 & 0.933 & 0.0008 & 0.944 \\
TQA Phi-3 & LEX+0D & -- & 0.821 & 0.0003 & 0.842 \\
TQA Phi-3 & LEX+SINK & -- & 0.820 & 0.0005 & 0.832 \\
TQA Phi-3 & PH SINK & 1024 & 0.871 & 0.0008 & 0.898 \\
TQA Phi-3 & PH SINK+PH 0DTOT & 2048 & 0.871 & 0.0007 & 0.906 \\
TQA Phi-3 & LEN & 2 & 0.522 & 0.0004 & 0.562 \\
TQA Phi-3 & 0D & 64 & 0.647 & 0.0010 & 0.717 \\
TQA TinyLl. & LEN & 2 & 0.523 & 0.0005 & 0.556 \\
TQA TinyLl. & 0D & 44 & 0.654 & 0.0012 & 0.727 \\
TQA TinyLl. & SINK & 44 & 0.644 & 0.0008 & 0.720 \\
TQA TinyLl. & 1D & 88 & 0.593 & 0.0010 & 0.677 \\
TQA TinyLl. & DEFL & 132 & 0.652 & 0.0028 & 0.704 \\
TQA TinyLl. & ANS & 132 & 0.646 & 0.0041 & 0.683 \\
TQA TinyLl. & ROWSTAT & 44 & 0.660 & 0.0016 & 0.744 \\
TQA TinyLl. & LOGPROB & 6 & 0.637 & 0.0010 & 0.646 \\
TQA TinyLl. & HIDDEN & 2048 & 0.883 & 0.0013 & 0.879 \\
TQA TinyLl. & LEX & -- & 0.820 & 0.0014 & 0.829 \\
TQA TinyLl. & LEGACY MSTPROXY & 22 & 0.627 & 0.0008 & 0.738 \\
TQA TinyLl. & LLMCHECK & 22 & 0.614 & 0.0007 & 0.701 \\
TQA TinyLl. & LOOKBACK & 704 & 0.850 & 0.0028 & 0.846 \\
TQA TinyLl. & PH 0D & 1408 & 0.805 & 0.0022 & 0.821 \\
TQA TinyLl. & LEX+DEFL & -- & 0.824 & 0.0011 & 0.837 \\
TQA TinyLl. & PH SINK+PH 0DTOT & 1408 & 0.811 & 0.0027 & 0.840 \\
TQA TinyLl. & LEX+0D & -- & 0.829 & 0.0011 & 0.837 \\
TQA TinyLl. & HIDDEN+0D & 2092 & 0.883 & 0.0014 & 0.879 \\
TQA TinyLl. & LOGPROB+0D & 50 & 0.682 & 0.0010 & 0.717 \\
TQA TinyLl. & NONTOPO+DEFL & 2956 & 0.884 & 0.0009 & 0.882 \\
TQA TinyLl. & NONTOPO+0D & 2868 & 0.885 & 0.0012 & 0.887 \\
TQA TinyLl. & NONTOPO & 2824 & 0.885 & 0.0011 & 0.887 \\
TQA TinyLl. & PH SINK & 704 & 0.807 & 0.0025 & 0.845 \\
TQA TinyLl. & PH 0DTOT & 704 & 0.804 & 0.0018 & 0.831 \\
TQA TinyLl. & LEX+SINK & -- & 0.831 & 0.0012 & 0.832 \\
TQA TinyLl. & SINK+ANS & 176 & 0.687 & 0.0033 & 0.722 \\
TQA TinyLl. & SINK+DEFL & 176 & 0.689 & 0.0015 & 0.750 \\
TQA TinyLl. & 0D+1D & 132 & 0.659 & 0.0017 & 0.734 \\
TQA TinyLl. & 0D+SINK & 88 & 0.663 & 0.0004 & 0.723 \\
TQA TinyLl. & PH ENT & 704 & 0.817 & 0.0009 & 0.838 \\
TQA TinyLl. & SINK+1D & 132 & 0.651 & 0.0023 & 0.741 \\
TQA SmolLM & LEX+DEFL & -- & 0.822 & 0.0005 & 0.842 \\
TQA SmolLM & LEN & 2 & 0.524 & 0.0003 & 0.554 \\
TQA SmolLM & NONTOPO+0D & 2942 & 0.899 & 0.0012 & 0.902 \\
TQA SmolLM & PH SINK+PH 0DTOT & 1536 & 0.825 & 0.0021 & 0.856 \\
TQA SmolLM & NONTOPO & 2894 & 0.899 & 0.0015 & 0.903 \\
TQA SmolLM & SINK+DEFL & 192 & 0.652 & 0.0027 & 0.674 \\
TQA SmolLM & NONTOPO+DEFL & 3038 & 0.899 & 0.0011 & 0.897 \\
TQA SmolLM & LOGPROB+0D & 54 & 0.614 & 0.0011 & 0.628 \\
TQA SmolLM & HIDDEN+0D & 2096 & 0.894 & 0.0021 & 0.891 \\
TQA SmolLM & LEX+0D & -- & 0.827 & 0.0009 & 0.840 \\
TQA SmolLM & LEX+SINK & -- & 0.819 & 0.0015 & 0.825 \\
TQA SmolLM & 0D+1D & 144 & 0.588 & 0.0018 & 0.646 \\
TQA SmolLM & 0D+SINK & 96 & 0.641 & 0.0027 & 0.670 \\
TQA SmolLM & PH ENT & 768 & 0.834 & 0.0007 & 0.843 \\
TQA SmolLM & PH SINK & 768 & 0.720 & 0.0009 & 0.776 \\
TQA SmolLM & SINK+1D & 144 & 0.632 & 0.0028 & 0.645 \\
TQA SmolLM & PH 0DTOT & 768 & 0.819 & 0.0018 & 0.846 \\
TQA SmolLM & LOOKBACK & 768 & 0.863 & 0.0017 & 0.865 \\
TQA SmolLM & LLMCHECK & 24 & 0.604 & 0.0006 & 0.685 \\
TQA SmolLM & LEGACY MSTPROXY & 24 & 0.585 & 0.0011 & 0.681 \\
TQA SmolLM & LEX & -- & 0.820 & 0.0014 & 0.829 \\
TQA SmolLM & HIDDEN & 2048 & 0.894 & 0.0022 & 0.891 \\
TQA SmolLM & LOGPROB & 6 & 0.589 & 0.0010 & 0.591 \\
TQA SmolLM & ROWSTAT & 48 & 0.673 & 0.0013 & 0.742 \\
TQA SmolLM & ANS & 144 & 0.629 & 0.0034 & 0.634 \\
TQA SmolLM & DEFL & 144 & 0.592 & 0.0024 & 0.656 \\
TQA SmolLM & 1D & 96 & 0.545 & 0.0006 & 0.569 \\
TQA SmolLM & SINK & 48 & 0.611 & 0.0020 & 0.632 \\
TQA SmolLM & 0D & 48 & 0.584 & 0.0028 & 0.655 \\
TQA SmolLM & PH 0D & 1536 & 0.798 & 0.0011 & 0.820 \\
TQA SmolLM & SINK+ANS & 192 & 0.646 & 0.0024 & 0.655 \\
HE Qw3B & LOGPROB+0D & 78 & 0.993 & 0.0003 & 0.990 \\
HE Qw3B & LEX+DEFL & -- & 0.970 & 0.0005 & 0.977 \\
HE Qw3B & SINK+1D & 216 & 0.842 & 0.0004 & 0.938 \\
HE Qw3B & SINK+ANS & 288 & 0.978 & 0.0007 & 0.977 \\
HE Qw3B & SINK+DEFL & 288 & 0.882 & 0.0026 & 0.961 \\
HE Qw3B & 0D+1D & 216 & 0.830 & 0.0023 & 0.944 \\
HE Qw3B & 0D+SINK & 144 & 0.858 & 0.0009 & 0.946 \\
HE Qw3B & LEGACY MSTPROXY & 36 & 0.733 & 0.0001 & 0.960 \\
HE Qw3B & LEX+0D & -- & 0.979 & 0.0001 & 0.983 \\
HE Qw3B & HIDDEN & 2048 & 0.999 & 0.0000 & 0.999 \\
HE Qw3B & LOGPROB & 6 & 0.993 & 0.0003 & 0.986 \\
HE Qw3B & ROWSTAT & 72 & 0.830 & 0.0003 & 0.975 \\
HE Qw3B & ANS & 216 & 0.974 & 0.0008 & 0.974 \\
HE Qw3B & DEFL & 216 & 0.832 & 0.0001 & 0.958 \\
HE Qw3B & 1D & 144 & 0.709 & 0.0010 & 0.892 \\
HE Qw3B & SINK & 72 & 0.827 & 0.0008 & 0.930 \\
HE Qw3B & 0D & 72 & 0.816 & 0.0009 & 0.929 \\
HE Qw3B & LEN & 2 & 0.962 & 0.0001 & 0.970 \\
HE Qw3B & LEX+SINK & -- & 0.980 & 0.0002 & 0.983 \\
HE Qw3B & HIDDEN+0D & 2120 & 0.999 & 0.0001 & 0.999 \\
HE Qw3B & LEX & -- & 0.981 & 0.0004 & 0.977 \\
HE Qw1.5B & LEN & 2 & 0.967 & 0.0002 & 0.976 \\
HE Qw1.5B & LEX+DEFL & -- & 0.956 & 0.0002 & 0.986 \\
HE Qw1.5B & LEX+0D & -- & 0.971 & 0.0004 & 0.986 \\
HE Qw1.5B & HIDDEN+0D & 1592 & 0.999 & 0.0001 & 1.000 \\
HE Qw1.5B & LOGPROB+0D & 62 & 0.993 & 0.0004 & 0.994 \\
HE Qw1.5B & NONTOPO+DEFL & 2130 & 0.999 & 0.0004 & 1.000 \\
HE Qw1.5B & NONTOPO+0D & 2018 & 0.999 & 0.0001 & 1.000 \\
HE Qw1.5B & NONTOPO & 1962 & 0.999 & 0.0003 & 1.000 \\
HE Qw1.5B & PH SINK+PH 0DTOT & 672 & 0.931 & 0.0024 & 0.998 \\
HE Qw1.5B & SINK+1D & 168 & 0.826 & 0.0035 & 0.940 \\
HE Qw1.5B & SINK+ANS & 224 & 0.970 & 0.0022 & 0.990 \\
HE Qw1.5B & SINK+DEFL & 224 & 0.872 & 0.0028 & 0.962 \\
HE Qw1.5B & 0D+1D & 168 & 0.815 & 0.0032 & 0.922 \\
HE Qw1.5B & 0D+SINK & 112 & 0.832 & 0.0015 & 0.924 \\
HE Qw1.5B & PH ENT & 336 & 0.933 & 0.0011 & 0.996 \\
HE Qw1.5B & PH SINK & 336 & 0.922 & 0.0007 & 0.998 \\
HE Qw1.5B & LEX+SINK & -- & 0.968 & 0.0005 & 0.984 \\
HE Qw1.5B & PH 0D & 672 & 0.949 & 0.0027 & 0.974 \\
HE Qw1.5B & 0D & 56 & 0.789 & 0.0014 & 0.918 \\
HE Qw1.5B & SINK & 56 & 0.820 & 0.0015 & 0.934 \\
HE Qw1.5B & 1D & 112 & 0.747 & 0.0006 & 0.878 \\
HE Qw1.5B & DEFL & 168 & 0.825 & 0.0021 & 0.952 \\
HE Qw1.5B & ANS & 168 & 0.969 & 0.0010 & 0.983 \\
HE Qw1.5B & ROWSTAT & 56 & 0.788 & 0.0017 & 0.984 \\
HE Qw1.5B & PH 0DTOT & 336 & 0.921 & 0.0007 & 0.996 \\
HE Qw1.5B & HIDDEN & 1536 & 0.999 & 0.0001 & 1.000 \\
HE Qw1.5B & LEX & -- & 0.979 & 0.0004 & 0.980 \\
HE Qw1.5B & LEGACY MSTPROXY & 28 & 0.719 & 0.0003 & 0.976 \\
HE Qw1.5B & LLMCHECK & 28 & 0.791 & 0.0005 & 0.978 \\
HE Qw1.5B & LOOKBACK & 336 & 0.997 & 0.0001 & 0.996 \\
HE Qw1.5B & LOGPROB & 6 & 0.993 & 0.0001 & 0.992 \\
TrQA Qw3B & LOOKBACK & 576 & 0.861 & 0.0013 & -- \\
TrQA Qw3B & LLMCHECK & 36 & 0.688 & 0.0012 & -- \\
TrQA Qw3B & LEGACY MSTPROXY & 36 & 0.691 & 0.0033 & -- \\
TrQA Qw3B & LEX & -- & 0.596 & 0.0039 & -- \\
TrQA Qw3B & HIDDEN & 2048 & 0.841 & 0.0042 & -- \\
TrQA Qw3B & LOGPROB & 6 & 0.839 & 0.0002 & -- \\
TrQA Qw3B & DEFL & 216 & 0.684 & 0.0043 & -- \\
TrQA Qw3B & ANS & 216 & 0.719 & 0.0020 & -- \\
TrQA Qw3B & 1D & 144 & 0.601 & 0.0031 & -- \\
TrQA Qw3B & SINK & 72 & 0.717 & 0.0038 & -- \\
TrQA Qw3B & 0D & 72 & 0.732 & 0.0038 & -- \\
TrQA Qw3B & PH 0D & 1152 & 0.807 & 0.0017 & -- \\
TrQA Qw3B & ROWSTAT & 72 & 0.724 & 0.0046 & -- \\
TrQA Qw3B & LEN & 2 & 0.569 & 0.0014 & -- \\
TrQA Qw3B & PH 0DTOT & 576 & 0.807 & 0.0012 & -- \\
TrQA Qw3B & PH ENT & 576 & 0.810 & 0.0007 & -- \\
TrQA Qw3B & PH SINK & 576 & 0.820 & 0.0021 & -- \\
TrQA Qw3B & LEX+SINK & -- & 0.724 & 0.0005 & -- \\
TrQA Qw3B & LEX+0D & -- & 0.733 & 0.0011 & -- \\
TrQA Qw3B & HIDDEN+0D & 2120 & 0.843 & 0.0040 & -- \\
TrQA Qw3B & LOGPROB+0D & 78 & 0.843 & 0.0022 & -- \\
TrQA Qw3B & NONTOPO+DEFL & 2954 & 0.870 & 0.0029 & -- \\
TrQA Qw3B & NONTOPO+0D & 2810 & 0.870 & 0.0027 & -- \\
TrQA Qw3B & LEX+DEFL & -- & 0.690 & 0.0031 & -- \\
TrQA Qw3B & PH SINK+PH 0DTOT & 1152 & 0.824 & 0.0017 & -- \\
TrQA Qw3B & SINK+1D & 216 & 0.718 & 0.0035 & -- \\
TrQA Qw3B & SINK+ANS & 288 & 0.752 & 0.0010 & -- \\
TrQA Qw3B & SINK+DEFL & 288 & 0.719 & 0.0038 & -- \\
TrQA Qw3B & 0D+1D & 216 & 0.729 & 0.0022 & -- \\
TrQA Qw3B & 0D+SINK & 144 & 0.738 & 0.0050 & -- \\
TrQA Qw3B & NONTOPO & 2738 & 0.870 & 0.0028 & -- \\
TrQA Qw1.5B & LOOKBACK & 336 & 0.824 & 0.0017 & -- \\
TrQA Qw1.5B & LLMCHECK & 28 & 0.676 & 0.0029 & -- \\
TrQA Qw1.5B & LEGACY MSTPROXY & 28 & 0.685 & 0.0014 & -- \\
TrQA Qw1.5B & LEX & -- & 0.609 & 0.0033 & -- \\
TrQA Qw1.5B & HIDDEN & 1536 & 0.855 & 0.0031 & -- \\
TrQA Qw1.5B & LOGPROB & 6 & 0.823 & 0.0006 & -- \\
TrQA Qw1.5B & PH 0D & 672 & 0.768 & 0.0013 & -- \\
TrQA Qw1.5B & ANS & 168 & 0.675 & 0.0035 & -- \\
TrQA Qw1.5B & 1D & 112 & 0.574 & 0.0024 & -- \\
TrQA Qw1.5B & SINK & 56 & 0.707 & 0.0000 & -- \\
TrQA Qw1.5B & 0D & 56 & 0.716 & 0.0004 & -- \\
TrQA Qw1.5B & LEN & 2 & 0.588 & 0.0012 & -- \\
TrQA Qw1.5B & ROWSTAT & 56 & 0.720 & 0.0008 & -- \\
TrQA Qw1.5B & DEFL & 168 & 0.655 & 0.0016 & -- \\
TrQA Qw1.5B & PH 0DTOT & 336 & 0.759 & 0.0051 & -- \\
TrQA Qw1.5B & PH ENT & 336 & 0.768 & 0.0025 & -- \\
TrQA Qw1.5B & PH SINK & 336 & 0.767 & 0.0034 & -- \\
TrQA Qw1.5B & LEX+DEFL & -- & 0.672 & 0.0047 & -- \\
TrQA Qw1.5B & LEX+0D & -- & 0.721 & 0.0031 & -- \\
TrQA Qw1.5B & HIDDEN+0D & 1592 & 0.855 & 0.0032 & -- \\
TrQA Qw1.5B & LOGPROB+0D & 62 & 0.828 & 0.0007 & -- \\
TrQA Qw1.5B & NONTOPO+DEFL & 2130 & 0.872 & 0.0024 & -- \\
TrQA Qw1.5B & NONTOPO+0D & 2018 & 0.873 & 0.0025 & -- \\
TrQA Qw1.5B & LEX+SINK & -- & 0.715 & 0.0036 & -- \\
TrQA Qw1.5B & PH SINK+PH 0DTOT & 672 & 0.772 & 0.0022 & -- \\
TrQA Qw1.5B & SINK+1D & 168 & 0.709 & 0.0028 & -- \\
TrQA Qw1.5B & SINK+ANS & 224 & 0.726 & 0.0059 & -- \\
TrQA Qw1.5B & SINK+DEFL & 224 & 0.719 & 0.0023 & -- \\
TrQA Qw1.5B & NONTOPO & 1962 & 0.873 & 0.0023 & -- \\
TrQA Qw1.5B & 0D+1D & 168 & 0.714 & 0.0015 & -- \\
TrQA Qw1.5B & 0D+SINK & 112 & 0.717 & 0.0006 & -- \\
TrQA Phi-3 & LOOKBACK & 1024 & 0.840 & 0.0016 & -- \\
TrQA Phi-3 & LLMCHECK & 32 & 0.676 & 0.0018 & -- \\
TrQA Phi-3 & LEGACY MSTPROXY & 32 & 0.669 & 0.0002 & -- \\
TrQA Phi-3 & LEX & -- & 0.635 & 0.0083 & -- \\
TrQA Phi-3 & HIDDEN & 3072 & 0.867 & 0.0035 & -- \\
TrQA Phi-3 & ROWSTAT & 64 & 0.728 & 0.0006 & -- \\
TrQA Phi-3 & PH 0D & 2048 & 0.807 & 0.0024 & -- \\
TrQA Phi-3 & ANS & 192 & 0.710 & 0.0037 & -- \\
TrQA Phi-3 & DEFL & 192 & 0.698 & 0.0021 & -- \\
TrQA Phi-3 & 1D & 128 & 0.546 & 0.0031 & -- \\
TrQA Phi-3 & SINK & 64 & 0.707 & 0.0018 & -- \\
TrQA Phi-3 & 0D & 64 & 0.710 & 0.0022 & -- \\
TrQA Phi-3 & LOGPROB & 6 & 0.877 & 0.0002 & -- \\
TrQA Phi-3 & PH 0DTOT & 1024 & 0.810 & 0.0016 & -- \\
TrQA Phi-3 & 0D+1D & 192 & 0.705 & 0.0011 & -- \\
TrQA Phi-3 & PH ENT & 1024 & 0.817 & 0.0021 & -- \\
TrQA Phi-3 & 0D+SINK & 128 & 0.713 & 0.0024 & -- \\
TrQA Phi-3 & SINK+DEFL & 256 & 0.736 & 0.0016 & -- \\
TrQA Phi-3 & SINK+ANS & 256 & 0.736 & 0.0046 & -- \\
TrQA Phi-3 & SINK+1D & 192 & 0.699 & 0.0012 & -- \\
TrQA Phi-3 & NONTOPO & 4198 & 0.883 & 0.0030 & -- \\
TrQA Phi-3 & NONTOPO+0D & 4262 & 0.883 & 0.0030 & -- \\
TrQA Phi-3 & NONTOPO+DEFL & 4390 & 0.885 & 0.0030 & -- \\
TrQA Phi-3 & LOGPROB+0D & 70 & 0.878 & 0.0013 & -- \\
TrQA Phi-3 & HIDDEN+0D & 3136 & 0.867 & 0.0037 & -- \\
TrQA Phi-3 & LEX+0D & -- & 0.721 & 0.0031 & -- \\
TrQA Phi-3 & LEX+SINK & -- & 0.718 & 0.0044 & -- \\
TrQA Phi-3 & LEX+DEFL & -- & 0.712 & 0.0022 & -- \\
TrQA Phi-3 & LEN & 2 & 0.581 & 0.0022 & -- \\
TrQA Phi-3 & PH SINK & 1024 & 0.826 & 0.0010 & -- \\
TrQA Phi-3 & PH SINK+PH 0DTOT & 2048 & 0.828 & 0.0009 & -- \\
TrQA TinyLl. & LEX+SINK & -- & 0.762 & 0.0021 & -- \\
TrQA TinyLl. & LEN & 2 & 0.631 & 0.0005 & -- \\
TrQA TinyLl. & LEX+0D & -- & 0.764 & 0.0010 & -- \\
TrQA TinyLl. & HIDDEN+0D & 2092 & 0.828 & 0.0027 & -- \\
TrQA TinyLl. & LOGPROB+0D & 50 & 0.762 & 0.0021 & -- \\
TrQA TinyLl. & NONTOPO+DEFL & 2956 & 0.842 & 0.0032 & -- \\
TrQA TinyLl. & NONTOPO+0D & 2868 & 0.842 & 0.0029 & -- \\
TrQA TinyLl. & NONTOPO & 2824 & 0.842 & 0.0031 & -- \\
TrQA TinyLl. & PH SINK+PH 0DTOT & 1408 & 0.813 & 0.0016 & -- \\
TrQA TinyLl. & SINK+1D & 132 & 0.719 & 0.0013 & -- \\
TrQA TinyLl. & SINK+DEFL & 176 & 0.734 & 0.0032 & -- \\
TrQA TinyLl. & 0D+1D & 132 & 0.718 & 0.0002 & -- \\
TrQA TinyLl. & 0D+SINK & 88 & 0.724 & 0.0015 & -- \\
TrQA TinyLl. & PH ENT & 704 & 0.807 & 0.0003 & -- \\
TrQA TinyLl. & PH SINK & 704 & 0.814 & 0.0016 & -- \\
TrQA TinyLl. & PH 0DTOT & 704 & 0.792 & 0.0018 & -- \\
TrQA TinyLl. & PH 0D & 1408 & 0.775 & 0.0027 & -- \\
TrQA TinyLl. & LOOKBACK & 704 & 0.823 & 0.0024 & -- \\
TrQA TinyLl. & LLMCHECK & 22 & 0.688 & 0.0010 & -- \\
TrQA TinyLl. & LEGACY MSTPROXY & 22 & 0.729 & 0.0006 & -- \\
TrQA TinyLl. & LEX & -- & 0.740 & 0.0003 & -- \\
TrQA TinyLl. & HIDDEN & 2048 & 0.827 & 0.0028 & -- \\
TrQA TinyLl. & LOGPROB & 6 & 0.724 & 0.0006 & -- \\
TrQA TinyLl. & ROWSTAT & 44 & 0.752 & 0.0001 & -- \\
TrQA TinyLl. & ANS & 132 & 0.701 & 0.0021 & -- \\
TrQA TinyLl. & DEFL & 132 & 0.726 & 0.0024 & -- \\
TrQA TinyLl. & 1D & 88 & 0.659 & 0.0007 & -- \\
TrQA TinyLl. & SINK & 44 & 0.720 & 0.0020 & -- \\
TrQA TinyLl. & 0D & 44 & 0.720 & 0.0020 & -- \\
TrQA TinyLl. & LEX+DEFL & -- & 0.757 & 0.0032 & -- \\
TrQA TinyLl. & SINK+ANS & 176 & 0.725 & 0.0030 & -- \\
\bottomrule
\end{longtable}

{\tiny\setlength{\tabcolsep}{2.5pt}\begin{longtable}{llrrrrcc}
\toprule
Setting & Test (A $\to$ B) & AUC$_A$ & AUC$_B$ & $\Delta$ & 90\% CI & $\equiv_{.015}$ & power \\
\midrule\endhead
TQA Qw3B & LEX$\to$LEX+0D & 0.820 & 0.847 & +0.027 & [+0.016, +0.038] & no & 0.48 \\
TQA Qw3B & 0D$\to$LEGACY\_MSTPROXY & 0.754 & 0.714 & -0.041 & [-0.053, -0.028] & no & 0.26 \\
TQA Qw3B & SINK$\to$0D+SINK & 0.740 & 0.757 & +0.017 & [+0.009, +0.026] & no & 0.81 \\
TQA Qw3B & SINK$\to$SINK+DEFL & 0.740 & 0.772 & +0.032 & [+0.019, +0.044] & no & 0.31 \\
TQA Qw3B & 0D$\to$0D+1D & 0.754 & 0.760 & +0.006 & [-0.002, +0.013] & yes & 0.89 \\
TQA Qw3B & SINK$\to$SINK+1D & 0.740 & 0.747 & +0.007 & [-0.001, +0.014] & yes & 0.90 \\
TQA Qw3B & SINK$\to$SINK+ANS & 0.740 & 0.784 & +0.044 & [+0.029, +0.059] & no & 0.04 \\
TQA Qw3B & PH\_0DTOT$\to$PH\_SINK & 0.844 & 0.850 & +0.007 & [-0.001, +0.014] & yes & 0.92 \\
TQA Qw3B & PH\_SINK$\to$PH\_SINK+PH\_0DTOT & 0.850 & 0.860 & +0.010 & [+0.006, +0.014] & yes & 1.00 \\
TQA Qw3B & PH\_0D$\to$PH\_ENT & 0.854 & 0.859 & +0.005 & [-0.005, +0.015] & no & 0.55 \\
TQA Qw3B & 0D$\to$SINK & 0.754 & 0.740 & -0.014 & [-0.024, -0.004] & no & 0.59 \\
TQA Qw3B & 0D$\to$LOGPROB & 0.754 & 0.631 & -0.124 & [-0.146, -0.101] & no & 0.00 \\
TQA Qw3B & 0D$\to$LOOKBACK & 0.754 & 0.897 & +0.143 & [+0.128, +0.157] & no & 0.03 \\
TQA Qw3B & 0D$\to$LLMCHECK & 0.754 & 0.693 & -0.061 & [-0.078, -0.044] & no & 0.00 \\
TQA Qw3B & 0D$\to$LEX & 0.754 & 0.820 & +0.066 & [+0.047, +0.084] & no & 0.00 \\
TQA Qw3B & 0D$\to$ROWSTAT & 0.754 & 0.751 & -0.003 & [-0.016, +0.010] & no & 0.22 \\
TQA Qw3B & 0D$\to$LEN & 0.754 & 0.529 & -0.226 & [-0.248, -0.203] & no & 0.00 \\
TQA Qw3B & HIDDEN$\to$HIDDEN+0D & 0.913 & 0.914 & +0.001 & [-0.000, +0.002] & yes & 1.00 \\
TQA Qw3B & LOGPROB$\to$LOGPROB+0D & 0.631 & 0.787 & +0.157 & [+0.138, +0.176] & no & 0.00 \\
TQA Qw3B & LEX$\to$LEX+DEFL & 0.820 & 0.829 & +0.009 & [-0.004, +0.022] & no & 0.22 \\
TQA Qw3B & LEX$\to$LEX+SINK & 0.820 & 0.843 & +0.023 & [+0.013, +0.033] & no & 0.57 \\
TQA Qw3B & 0D$\to$HIDDEN & 0.754 & 0.913 & +0.159 & [+0.144, +0.174] & no & 0.03 \\
TQA Qw3B & NONTOPO$\to$NONTOPO+0D & 0.927 & 0.928 & +0.001 & [-0.000, +0.001] & yes & 1.00 \\
TQA Qw3B & NONTOPO$\to$NONTOPO+DEFL & 0.927 & 0.928 & +0.001 & [-0.000, +0.003] & yes & 1.00 \\
TQA Qw1.5B & 0D$\to$LEGACY\_MSTPROXY & 0.726 & 0.698 & -0.028 & [-0.040, -0.016] & no & 0.36 \\
TQA Qw1.5B & 0D$\to$SINK & 0.726 & 0.723 & -0.003 & [-0.012, +0.006] & yes & 0.68 \\
TQA Qw1.5B & LOGPROB$\to$LOGPROB+0D & 0.638 & 0.740 & +0.101 & [+0.086, +0.116] & no & 0.00 \\
TQA Qw1.5B & LEX$\to$LEX+DEFL & 0.820 & 0.825 & +0.005 & [-0.006, +0.016] & no & 0.45 \\
TQA Qw1.5B & LEX$\to$LEX+SINK & 0.820 & 0.833 & +0.013 & [+0.004, +0.022] & no & 0.73 \\
TQA Qw1.5B & LEX$\to$LEX+0D & 0.820 & 0.835 & +0.015 & [+0.006, +0.024] & no & 0.71 \\
TQA Qw1.5B & NONTOPO$\to$NONTOPO+DEFL & 0.917 & 0.917 & +0.000 & [-0.001, +0.001] & yes & 1.00 \\
TQA Qw1.5B & NONTOPO$\to$NONTOPO+0D & 0.917 & 0.917 & +0.000 & [-0.000, +0.001] & yes & 1.00 \\
TQA Qw1.5B & 0D$\to$LEN & 0.726 & 0.529 & -0.197 & [-0.218, -0.176] & no & 0.00 \\
TQA Qw1.5B & 0D$\to$ROWSTAT & 0.726 & 0.738 & +0.012 & [-0.000, +0.025] & no & 0.27 \\
TQA Qw1.5B & 0D$\to$LEX & 0.726 & 0.820 & +0.094 & [+0.076, +0.112] & no & 0.00 \\
TQA Qw1.5B & 0D$\to$LLMCHECK & 0.726 & 0.666 & -0.060 & [-0.076, -0.044] & no & 0.00 \\
TQA Qw1.5B & HIDDEN$\to$HIDDEN+0D & 0.916 & 0.916 & +0.001 & [-0.000, +0.002] & yes & 1.00 \\
TQA Qw1.5B & 0D$\to$HIDDEN & 0.726 & 0.916 & +0.190 & [+0.175, +0.204] & no & 0.05 \\
TQA Qw1.5B & 0D$\to$LOGPROB & 0.726 & 0.638 & -0.088 & [-0.108, -0.067] & no & 0.00 \\
TQA Qw1.5B & PH\_0D$\to$PH\_ENT & 0.838 & 0.833 & -0.004 & [-0.016, +0.007] & no & 0.44 \\
TQA Qw1.5B & PH\_SINK$\to$PH\_SINK+PH\_0DTOT & 0.830 & 0.841 & +0.011 & [+0.007, +0.014] & yes & 1.00 \\
TQA Qw1.5B & PH\_0DTOT$\to$PH\_SINK & 0.819 & 0.830 & +0.011 & [+0.003, +0.018] & no & 0.90 \\
TQA Qw1.5B & SINK$\to$SINK+ANS & 0.723 & 0.761 & +0.038 & [+0.024, +0.051] & no & 0.11 \\
TQA Qw1.5B & SINK$\to$SINK+1D & 0.723 & 0.715 & -0.008 & [-0.014, -0.002] & yes & 0.99 \\
TQA Qw1.5B & 0D$\to$0D+1D & 0.726 & 0.718 & -0.008 & [-0.014, -0.002] & yes & 0.99 \\
TQA Qw1.5B & SINK$\to$SINK+DEFL & 0.723 & 0.757 & +0.034 & [+0.022, +0.045] & no & 0.35 \\
TQA Qw1.5B & SINK$\to$0D+SINK & 0.723 & 0.735 & +0.012 & [+0.005, +0.020] & no & 0.91 \\
TQA Qw1.5B & 0D$\to$LOOKBACK & 0.726 & 0.851 & +0.125 & [+0.109, +0.140] & no & 0.00 \\
TQA Phi-3 & PH\_0D$\to$PH\_ENT & 0.868 & 0.864 & -0.004 & [-0.012, +0.004] & yes & 0.87 \\
TQA Phi-3 & PH\_SINK$\to$PH\_SINK+PH\_0DTOT & 0.871 & 0.871 & +0.000 & [-0.001, +0.001] & yes & 1.00 \\
TQA Phi-3 & PH\_0DTOT$\to$PH\_SINK & 0.846 & 0.871 & +0.025 & [+0.020, +0.030] & no & 1.00 \\
TQA Phi-3 & SINK$\to$0D+SINK & 0.652 & 0.658 & +0.006 & [+0.000, +0.012] & yes & 0.99 \\
TQA Phi-3 & SINK$\to$SINK+1D & 0.652 & 0.654 & +0.002 & [-0.003, +0.006] & yes & 1.00 \\
TQA Phi-3 & 0D$\to$0D+1D & 0.647 & 0.645 & -0.001 & [-0.005, +0.002] & yes & 1.00 \\
TQA Phi-3 & SINK$\to$SINK+DEFL & 0.652 & 0.685 & +0.033 & [+0.019, +0.047] & no & 0.08 \\
TQA Phi-3 & 0D$\to$LOGPROB & 0.647 & 0.578 & -0.069 & [-0.093, -0.045] & no & 0.00 \\
TQA Phi-3 & SINK$\to$SINK+ANS & 0.652 & 0.672 & +0.020 & [+0.004, +0.036] & no & 0.00 \\
TQA Phi-3 & 0D$\to$HIDDEN & 0.647 & 0.933 & +0.286 & [+0.271, +0.301] & no & 0.00 \\
TQA Phi-3 & HIDDEN$\to$HIDDEN+0D & 0.933 & 0.933 & +0.000 & [-0.000, +0.001] & yes & 1.00 \\
TQA Phi-3 & 0D$\to$LLMCHECK & 0.647 & 0.653 & +0.007 & [-0.011, +0.024] & no & 0.00 \\
TQA Phi-3 & 0D$\to$LEX & 0.647 & 0.820 & +0.173 & [+0.153, +0.193] & no & 0.00 \\
TQA Phi-3 & 0D$\to$ROWSTAT & 0.647 & 0.661 & +0.014 & [+0.001, +0.028] & no & 0.12 \\
TQA Phi-3 & NONTOPO$\to$NONTOPO+0D & 0.936 & 0.936 & +0.000 & [-0.000, +0.000] & yes & 1.00 \\
TQA Phi-3 & NONTOPO$\to$NONTOPO+DEFL & 0.936 & 0.937 & +0.001 & [-0.000, +0.002] & yes & 1.00 \\
TQA Phi-3 & LEX$\to$LEX+0D & 0.820 & 0.821 & +0.001 & [-0.007, +0.009] & yes & 0.85 \\
TQA Phi-3 & LEX$\to$LEX+SINK & 0.820 & 0.820 & -0.000 & [-0.008, +0.008] & yes & 0.88 \\
TQA Phi-3 & LEX$\to$LEX+DEFL & 0.820 & 0.802 & -0.018 & [-0.028, -0.007] & no & 0.52 \\
TQA Phi-3 & LOGPROB$\to$LOGPROB+0D & 0.578 & 0.656 & +0.078 & [+0.059, +0.097] & no & 0.00 \\
TQA Phi-3 & 0D$\to$LOOKBACK & 0.647 & 0.892 & +0.245 & [+0.229, +0.261] & no & 0.00 \\
TQA Phi-3 & 0D$\to$LEN & 0.647 & 0.522 & -0.124 & [-0.146, -0.103] & no & 0.00 \\
TQA Phi-3 & 0D$\to$SINK & 0.647 & 0.652 & +0.005 & [-0.002, +0.013] & yes & 0.91 \\
TQA Phi-3 & 0D$\to$LEGACY\_MSTPROXY & 0.647 & 0.614 & -0.033 & [-0.046, -0.020] & no & 0.20 \\
TQA TinyLl. & 0D$\to$SINK & 0.654 & 0.644 & -0.010 & [-0.021, +0.001] & no & 0.48 \\
TQA TinyLl. & 0D$\to$LEGACY\_MSTPROXY & 0.654 & 0.627 & -0.027 & [-0.038, -0.016] & no & 0.45 \\
TQA TinyLl. & SINK$\to$0D+SINK & 0.644 & 0.663 & +0.020 & [+0.010, +0.029] & no & 0.72 \\
TQA TinyLl. & SINK$\to$SINK+DEFL & 0.644 & 0.689 & +0.045 & [+0.033, +0.058] & no & 0.25 \\
TQA TinyLl. & 0D$\to$0D+1D & 0.654 & 0.659 & +0.005 & [-0.002, +0.012] & yes & 0.94 \\
TQA TinyLl. & SINK$\to$SINK+1D & 0.644 & 0.651 & +0.007 & [-0.000, +0.015] & yes & 0.90 \\
TQA TinyLl. & SINK$\to$SINK+ANS & 0.644 & 0.687 & +0.043 & [+0.026, +0.060] & no & 0.00 \\
TQA TinyLl. & PH\_0DTOT$\to$PH\_SINK & 0.804 & 0.807 & +0.003 & [-0.003, +0.009] & yes & 0.97 \\
TQA TinyLl. & PH\_SINK$\to$PH\_SINK+PH\_0DTOT & 0.807 & 0.811 & +0.004 & [+0.002, +0.005] & yes & 1.00 \\
TQA TinyLl. & PH\_0D$\to$PH\_ENT & 0.805 & 0.817 & +0.011 & [+0.001, +0.022] & no & 0.50 \\
TQA TinyLl. & 0D$\to$LOGPROB & 0.654 & 0.637 & -0.017 & [-0.039, +0.004] & no & 0.00 \\
TQA TinyLl. & HIDDEN$\to$HIDDEN+0D & 0.883 & 0.883 & -0.000 & [-0.001, +0.000] & yes & 1.00 \\
TQA TinyLl. & 0D$\to$LOOKBACK & 0.654 & 0.850 & +0.196 & [+0.182, +0.212] & no & 0.00 \\
TQA TinyLl. & 0D$\to$LLMCHECK & 0.654 & 0.614 & -0.040 & [-0.055, -0.024] & no & 0.00 \\
TQA TinyLl. & 0D$\to$LEX & 0.654 & 0.820 & +0.166 & [+0.149, +0.184] & no & 0.00 \\
TQA TinyLl. & 0D$\to$ROWSTAT & 0.654 & 0.660 & +0.006 & [-0.005, +0.018] & no & 0.37 \\
TQA TinyLl. & LOGPROB$\to$LOGPROB+0D & 0.637 & 0.682 & +0.046 & [+0.032, +0.059] & no & 0.13 \\
TQA TinyLl. & 0D$\to$LEN & 0.654 & 0.523 & -0.131 & [-0.152, -0.110] & no & 0.00 \\
TQA TinyLl. & NONTOPO$\to$NONTOPO+0D & 0.885 & 0.885 & +0.000 & [-0.000, +0.000] & yes & 1.00 \\
TQA TinyLl. & NONTOPO$\to$NONTOPO+DEFL & 0.885 & 0.884 & -0.001 & [-0.002, -0.000] & yes & 1.00 \\
TQA TinyLl. & LEX$\to$LEX+0D & 0.820 & 0.829 & +0.009 & [+0.003, +0.016] & no & 0.98 \\
TQA TinyLl. & LEX$\to$LEX+SINK & 0.820 & 0.831 & +0.011 & [+0.005, +0.017] & no & 0.98 \\
TQA TinyLl. & LEX$\to$LEX+DEFL & 0.820 & 0.824 & +0.004 & [-0.005, +0.013] & yes & 0.75 \\
TQA TinyLl. & 0D$\to$HIDDEN & 0.654 & 0.883 & +0.229 & [+0.213, +0.244] & no & 0.00 \\
TQA SmolLM & HIDDEN$\to$HIDDEN+0D & 0.894 & 0.894 & +0.000 & [-0.001, +0.001] & yes & 1.00 \\
TQA SmolLM & 0D$\to$SINK & 0.584 & 0.611 & +0.027 & [+0.008, +0.046] & no & 0.00 \\
TQA SmolLM & NONTOPO$\to$NONTOPO+DEFL & 0.899 & 0.899 & -0.000 & [-0.001, +0.001] & yes & 1.00 \\
TQA SmolLM & 0D$\to$ROWSTAT & 0.584 & 0.673 & +0.090 & [+0.078, +0.101] & no & 0.36 \\
TQA SmolLM & LEX$\to$LEX+0D & 0.820 & 0.827 & +0.007 & [+0.001, +0.012] & yes & 0.99 \\
TQA SmolLM & LEX$\to$LEX+SINK & 0.820 & 0.819 & -0.001 & [-0.007, +0.004] & yes & 1.00 \\
TQA SmolLM & LEX$\to$LEX+DEFL & 0.820 & 0.822 & +0.002 & [-0.006, +0.009] & yes & 0.90 \\
TQA SmolLM & LOGPROB$\to$LOGPROB+0D & 0.589 & 0.614 & +0.025 & [+0.015, +0.035] & no & 0.58 \\
TQA SmolLM & 0D$\to$LEX & 0.584 & 0.820 & +0.236 & [+0.217, +0.255] & no & 0.00 \\
TQA SmolLM & 0D$\to$LLMCHECK & 0.584 & 0.604 & +0.020 & [+0.005, +0.035] & no & 0.00 \\
TQA SmolLM & 0D$\to$LOOKBACK & 0.584 & 0.863 & +0.279 & [+0.264, +0.295] & no & 0.00 \\
TQA SmolLM & 0D$\to$HIDDEN & 0.584 & 0.894 & +0.310 & [+0.295, +0.325] & no & 0.02 \\
TQA SmolLM & 0D$\to$LOGPROB & 0.584 & 0.589 & +0.005 & [-0.017, +0.027] & no & 0.00 \\
TQA SmolLM & PH\_0D$\to$PH\_ENT & 0.798 & 0.834 & +0.036 & [+0.026, +0.046] & no & 0.59 \\
TQA SmolLM & PH\_SINK$\to$PH\_SINK+PH\_0DTOT & 0.720 & 0.825 & +0.105 & [+0.093, +0.118] & no & 0.29 \\
TQA SmolLM & PH\_0DTOT$\to$PH\_SINK & 0.819 & 0.720 & -0.099 & [-0.113, -0.086] & no & 0.16 \\
TQA SmolLM & SINK$\to$SINK+ANS & 0.611 & 0.646 & +0.036 & [+0.018, +0.053] & no & 0.00 \\
TQA SmolLM & SINK$\to$SINK+1D & 0.611 & 0.632 & +0.022 & [+0.012, +0.032] & no & 0.58 \\
TQA SmolLM & 0D$\to$0D+1D & 0.584 & 0.588 & +0.004 & [-0.005, +0.014] & yes & 0.66 \\
TQA SmolLM & SINK$\to$SINK+DEFL & 0.611 & 0.652 & +0.041 & [+0.030, +0.053] & no & 0.36 \\
TQA SmolLM & SINK$\to$0D+SINK & 0.611 & 0.641 & +0.030 & [+0.021, +0.039] & no & 0.67 \\
TQA SmolLM & 0D$\to$LEGACY\_MSTPROXY & 0.584 & 0.585 & +0.001 & [-0.004, +0.005] & yes & 1.00 \\
TQA SmolLM & NONTOPO$\to$NONTOPO+0D & 0.899 & 0.899 & -0.000 & [-0.001, +0.000] & yes & 1.00 \\
TQA SmolLM & 0D$\to$LEN & 0.584 & 0.524 & -0.060 & [-0.075, -0.044] & no & 0.00 \\
HE Qw3B & LEX$\to$LEX+0D & 0.981 & 0.979 & -0.002 & [-0.005, +0.000] & yes & 1.00 \\
HE Qw3B & HIDDEN$\to$HIDDEN+0D & 0.999 & 0.999 & -0.000 & [-0.000, +0.000] & yes & 1.00 \\
HE Qw3B & LOGPROB$\to$LOGPROB+0D & 0.993 & 0.993 & -0.000 & [-0.002, +0.001] & yes & 1.00 \\
HE Qw3B & 0D$\to$SINK & 0.816 & 0.827 & +0.011 & [+0.002, +0.021] & no & 0.66 \\
HE Qw3B & 0D$\to$LEGACY\_MSTPROXY & 0.816 & 0.733 & -0.083 & [-0.094, -0.072] & no & 0.45 \\
HE Qw3B & SINK$\to$0D+SINK & 0.827 & 0.858 & +0.030 & [+0.023, +0.037] & no & 0.96 \\
HE Qw3B & SINK$\to$SINK+DEFL & 0.827 & 0.882 & +0.055 & [+0.046, +0.063] & no & 0.75 \\
HE Qw3B & 0D$\to$0D+1D & 0.816 & 0.830 & +0.013 & [+0.007, +0.020] & no & 0.96 \\
HE Qw3B & SINK$\to$SINK+1D & 0.827 & 0.842 & +0.014 & [+0.008, +0.021] & no & 0.96 \\
HE Qw3B & SINK$\to$SINK+ANS & 0.827 & 0.978 & +0.151 & [+0.139, +0.162] & no & 0.41 \\
HE Qw3B & LEX$\to$LEX+DEFL & 0.981 & 0.970 & -0.011 & [-0.015, -0.007] & yes & 1.00 \\
HE Qw3B & 0D$\to$HIDDEN & 0.816 & 0.999 & +0.183 & [+0.171, +0.195] & no & 0.33 \\
HE Qw3B & 0D$\to$LEX & 0.816 & 0.981 & +0.165 & [+0.153, +0.176] & no & 0.38 \\
HE Qw3B & 0D$\to$ROWSTAT & 0.816 & 0.830 & +0.014 & [+0.002, +0.025] & no & 0.34 \\
HE Qw3B & 0D$\to$LEN & 0.816 & 0.962 & +0.146 & [+0.134, +0.158] & no & 0.37 \\
HE Qw3B & LEX$\to$LEX+SINK & 0.981 & 0.980 & -0.001 & [-0.004, +0.001] & yes & 1.00 \\
HE Qw3B & 0D$\to$LOGPROB & 0.816 & 0.993 & +0.177 & [+0.165, +0.189] & no & 0.35 \\
HE Qw1.5B & 0D$\to$SINK & 0.789 & 0.820 & +0.031 & [+0.017, +0.045] & no & 0.08 \\
HE Qw1.5B & HIDDEN$\to$HIDDEN+0D & 0.999 & 0.999 & +0.000 & [-0.000, +0.000] & yes & 1.00 \\
HE Qw1.5B & LOGPROB$\to$LOGPROB+0D & 0.993 & 0.993 & -0.001 & [-0.003, +0.001] & yes & 1.00 \\
HE Qw1.5B & LEX$\to$LEX+DEFL & 0.979 & 0.956 & -0.022 & [-0.030, -0.016] & no & 0.94 \\
HE Qw1.5B & LEX$\to$LEX+SINK & 0.979 & 0.968 & -0.011 & [-0.017, -0.005] & no & 0.98 \\
HE Qw1.5B & LEX$\to$LEX+0D & 0.979 & 0.971 & -0.008 & [-0.014, -0.002] & yes & 0.99 \\
HE Qw1.5B & NONTOPO$\to$NONTOPO+DEFL & 0.999 & 0.999 & +0.000 & [-0.000, +0.000] & yes & 1.00 \\
HE Qw1.5B & NONTOPO$\to$NONTOPO+0D & 0.999 & 0.999 & +0.000 & [+0.000, +0.001] & yes & 1.00 \\
HE Qw1.5B & 0D$\to$ROWSTAT & 0.789 & 0.788 & -0.001 & [-0.018, +0.016] & no & 0.00 \\
HE Qw1.5B & 0D$\to$LEX & 0.789 & 0.979 & +0.190 & [+0.172, +0.207] & no & 0.00 \\
HE Qw1.5B & 0D$\to$LLMCHECK & 0.789 & 0.791 & +0.002 & [-0.017, +0.021] & no & 0.00 \\
HE Qw1.5B & 0D$\to$LOOKBACK & 0.789 & 0.997 & +0.208 & [+0.191, +0.225] & no & 0.00 \\
HE Qw1.5B & 0D$\to$LEN & 0.789 & 0.967 & +0.178 & [+0.160, +0.196] & no & 0.00 \\
HE Qw1.5B & 0D$\to$LOGPROB & 0.789 & 0.993 & +0.204 & [+0.187, +0.221] & no & 0.00 \\
HE Qw1.5B & PH\_0D$\to$PH\_ENT & 0.949 & 0.933 & -0.016 & [-0.028, -0.004] & no & 0.33 \\
HE Qw1.5B & PH\_SINK$\to$PH\_SINK+PH\_0DTOT & 0.922 & 0.931 & +0.008 & [+0.006, +0.011] & yes & 1.00 \\
HE Qw1.5B & PH\_0DTOT$\to$PH\_SINK & 0.921 & 0.922 & +0.002 & [-0.009, +0.013] & yes & 0.46 \\
HE Qw1.5B & SINK$\to$SINK+ANS & 0.820 & 0.970 & +0.151 & [+0.133, +0.168] & no & 0.00 \\
HE Qw1.5B & SINK$\to$SINK+1D & 0.820 & 0.826 & +0.006 & [-0.005, +0.017] & no & 0.45 \\
HE Qw1.5B & 0D$\to$0D+1D & 0.789 & 0.815 & +0.026 & [+0.013, +0.040] & no & 0.17 \\
HE Qw1.5B & SINK$\to$SINK+DEFL & 0.820 & 0.872 & +0.052 & [+0.038, +0.065] & no & 0.17 \\
HE Qw1.5B & SINK$\to$0D+SINK & 0.820 & 0.832 & +0.012 & [+0.005, +0.019] & no & 0.91 \\
HE Qw1.5B & 0D$\to$HIDDEN & 0.789 & 0.999 & +0.210 & [+0.193, +0.227] & no & 0.00 \\
HE Qw1.5B & 0D$\to$LEGACY\_MSTPROXY & 0.789 & 0.719 & -0.070 & [-0.086, -0.053] & no & 0.00 \\
TrQA Qw3B & PH\_SINK$\to$PH\_SINK+PH\_0DTOT & 0.820 & 0.824 & +0.004 & [+0.002, +0.006] & yes & 1.00 \\
TrQA Qw3B & PH\_0DTOT$\to$PH\_SINK & 0.807 & 0.820 & +0.012 & [+0.004, +0.021] & no & 0.82 \\
TrQA Qw3B & SINK$\to$SINK+ANS & 0.717 & 0.752 & +0.035 & [+0.019, +0.051] & no & 0.00 \\
TrQA Qw3B & SINK$\to$SINK+1D & 0.717 & 0.718 & +0.001 & [-0.009, +0.010] & yes & 0.65 \\
TrQA Qw3B & 0D$\to$SINK & 0.732 & 0.717 & -0.015 & [-0.026, -0.004] & no & 0.48 \\
TrQA Qw3B & SINK$\to$SINK+DEFL & 0.717 & 0.719 & +0.002 & [-0.012, +0.015] & no & 0.13 \\
TrQA Qw3B & SINK$\to$0D+SINK & 0.717 & 0.738 & +0.021 & [+0.012, +0.030] & no & 0.76 \\
TrQA Qw3B & 0D$\to$LEGACY\_MSTPROXY & 0.732 & 0.691 & -0.041 & [-0.057, -0.025] & no & 0.00 \\
TrQA Qw3B & PH\_0D$\to$PH\_ENT & 0.807 & 0.810 & +0.003 & [-0.011, +0.017] & no & 0.08 \\
TrQA Qw3B & 0D$\to$0D+1D & 0.732 & 0.729 & -0.003 & [-0.012, +0.006] & yes & 0.74 \\
TrQA Qw3B & 0D$\to$LOGPROB & 0.732 & 0.839 & +0.107 & [+0.088, +0.125] & no & 0.00 \\
TrQA Qw3B & NONTOPO$\to$NONTOPO+0D & 0.870 & 0.870 & +0.001 & [-0.000, +0.002] & yes & 1.00 \\
TrQA Qw3B & 0D$\to$LOOKBACK & 0.732 & 0.861 & +0.129 & [+0.111, +0.147] & no & 0.00 \\
TrQA Qw3B & 0D$\to$LLMCHECK & 0.732 & 0.688 & -0.044 & [-0.065, -0.023] & no & 0.00 \\
TrQA Qw3B & 0D$\to$LEX & 0.732 & 0.596 & -0.137 & [-0.163, -0.110] & no & 0.00 \\
TrQA Qw3B & 0D$\to$ROWSTAT & 0.732 & 0.724 & -0.009 & [-0.025, +0.008] & no & 0.00 \\
TrQA Qw3B & 0D$\to$LEN & 0.732 & 0.569 & -0.163 & [-0.189, -0.137] & no & 0.00 \\
TrQA Qw3B & NONTOPO$\to$NONTOPO+DEFL & 0.870 & 0.870 & +0.000 & [-0.001, +0.002] & yes & 1.00 \\
TrQA Qw3B & LEX$\to$LEX+0D & 0.596 & 0.733 & +0.137 & [+0.117, +0.158] & no & 0.00 \\
TrQA Qw3B & LEX$\to$LEX+SINK & 0.596 & 0.724 & +0.128 & [+0.108, +0.149] & no & 0.00 \\
TrQA Qw3B & LEX$\to$LEX+DEFL & 0.596 & 0.690 & +0.095 & [+0.073, +0.116] & no & 0.00 \\
TrQA Qw3B & LOGPROB$\to$LOGPROB+0D & 0.839 & 0.843 & +0.004 & [-0.002, +0.011] & yes & 0.95 \\
TrQA Qw3B & HIDDEN$\to$HIDDEN+0D & 0.841 & 0.843 & +0.002 & [+0.000, +0.004] & yes & 1.00 \\
TrQA Qw3B & 0D$\to$HIDDEN & 0.732 & 0.841 & +0.109 & [+0.090, +0.129] & no & 0.00 \\
TrQA Qw1.5B & PH\_0D$\to$PH\_ENT & 0.768 & 0.768 & -0.000 & [-0.017, +0.017] & no & 0.00 \\
TrQA Qw1.5B & PH\_SINK$\to$PH\_SINK+PH\_0DTOT & 0.767 & 0.772 & +0.004 & [+0.001, +0.008] & yes & 1.00 \\
TrQA Qw1.5B & PH\_0DTOT$\to$PH\_SINK & 0.759 & 0.767 & +0.008 & [-0.003, +0.018] & no & 0.53 \\
TrQA Qw1.5B & SINK$\to$SINK+ANS & 0.707 & 0.726 & +0.019 & [+0.001, +0.036] & no & 0.00 \\
TrQA Qw1.5B & 0D$\to$SINK & 0.716 & 0.707 & -0.009 & [-0.020, +0.002] & no & 0.43 \\
TrQA Qw1.5B & 0D$\to$0D+1D & 0.716 & 0.714 & -0.002 & [-0.011, +0.008] & yes & 0.65 \\
TrQA Qw1.5B & SINK$\to$SINK+DEFL & 0.707 & 0.719 & +0.012 & [-0.003, +0.027] & no & 0.00 \\
TrQA Qw1.5B & SINK$\to$0D+SINK & 0.707 & 0.717 & +0.010 & [+0.001, +0.020] & no & 0.68 \\
TrQA Qw1.5B & 0D$\to$LEGACY\_MSTPROXY & 0.716 & 0.685 & -0.031 & [-0.046, -0.015] & no & 0.00 \\
TrQA Qw1.5B & SINK$\to$SINK+1D & 0.707 & 0.709 & +0.002 & [-0.008, +0.013] & yes & 0.49 \\
TrQA Qw1.5B & 0D$\to$HIDDEN & 0.716 & 0.855 & +0.140 & [+0.119, +0.160] & no & 0.00 \\
TrQA Qw1.5B & 0D$\to$LOGPROB & 0.716 & 0.823 & +0.107 & [+0.085, +0.129] & no & 0.00 \\
TrQA Qw1.5B & 0D$\to$LLMCHECK & 0.716 & 0.676 & -0.039 & [-0.062, -0.016] & no & 0.00 \\
TrQA Qw1.5B & 0D$\to$LOOKBACK & 0.716 & 0.824 & +0.108 & [+0.088, +0.129] & no & 0.00 \\
TrQA Qw1.5B & HIDDEN$\to$HIDDEN+0D & 0.855 & 0.855 & +0.000 & [-0.001, +0.002] & yes & 1.00 \\
TrQA Qw1.5B & LEX$\to$LEX+DEFL & 0.609 & 0.672 & +0.063 & [+0.040, +0.085] & no & 0.00 \\
TrQA Qw1.5B & LEX$\to$LEX+SINK & 0.609 & 0.715 & +0.106 & [+0.085, +0.127] & no & 0.00 \\
TrQA Qw1.5B & LEX$\to$LEX+0D & 0.609 & 0.721 & +0.112 & [+0.091, +0.134] & no & 0.00 \\
TrQA Qw1.5B & LOGPROB$\to$LOGPROB+0D & 0.823 & 0.828 & +0.006 & [-0.001, +0.013] & yes & 0.92 \\
TrQA Qw1.5B & NONTOPO$\to$NONTOPO+0D & 0.873 & 0.873 & +0.000 & [-0.001, +0.001] & yes & 1.00 \\
TrQA Qw1.5B & 0D$\to$LEN & 0.716 & 0.588 & -0.128 & [-0.154, -0.102] & no & 0.00 \\
TrQA Qw1.5B & 0D$\to$ROWSTAT & 0.716 & 0.720 & +0.005 & [-0.011, +0.020] & no & 0.00 \\
TrQA Qw1.5B & 0D$\to$LEX & 0.716 & 0.609 & -0.106 & [-0.136, -0.078] & no & 0.00 \\
TrQA Qw1.5B & NONTOPO$\to$NONTOPO+DEFL & 0.873 & 0.872 & -0.001 & [-0.002, +0.001] & yes & 1.00 \\
TrQA Phi-3 & 0D$\to$LEGACY\_MSTPROXY & 0.710 & 0.669 & -0.041 & [-0.057, -0.025] & no & 0.00 \\
TrQA Phi-3 & SINK$\to$0D+SINK & 0.707 & 0.713 & +0.006 & [+0.000, +0.012] & yes & 0.99 \\
TrQA Phi-3 & SINK$\to$SINK+DEFL & 0.707 & 0.736 & +0.029 & [+0.013, +0.045] & no & 0.00 \\
TrQA Phi-3 & 0D$\to$0D+1D & 0.710 & 0.705 & -0.005 & [-0.011, +0.000] & yes & 0.99 \\
TrQA Phi-3 & SINK$\to$SINK+1D & 0.707 & 0.699 & -0.008 & [-0.013, -0.002] & yes & 1.00 \\
TrQA Phi-3 & SINK$\to$SINK+ANS & 0.707 & 0.736 & +0.029 & [+0.012, +0.046] & no & 0.00 \\
TrQA Phi-3 & PH\_0DTOT$\to$PH\_SINK & 0.810 & 0.826 & +0.016 & [+0.009, +0.022] & no & 0.96 \\
TrQA Phi-3 & PH\_SINK$\to$PH\_SINK+PH\_0DTOT & 0.826 & 0.828 & +0.002 & [-0.001, +0.004] & yes & 1.00 \\
TrQA Phi-3 & PH\_0D$\to$PH\_ENT & 0.807 & 0.817 & +0.010 & [-0.002, +0.022] & no & 0.31 \\
TrQA Phi-3 & 0D$\to$LOGPROB & 0.710 & 0.877 & +0.167 & [+0.148, +0.186] & no & 0.00 \\
TrQA Phi-3 & 0D$\to$HIDDEN & 0.710 & 0.867 & +0.157 & [+0.137, +0.176] & no & 0.00 \\
TrQA Phi-3 & 0D$\to$LLMCHECK & 0.710 & 0.676 & -0.034 & [-0.055, -0.014] & no & 0.00 \\
TrQA Phi-3 & 0D$\to$LEX & 0.710 & 0.635 & -0.076 & [-0.102, -0.050] & no & 0.00 \\
TrQA Phi-3 & 0D$\to$ROWSTAT & 0.710 & 0.728 & +0.018 & [+0.002, +0.034] & no & 0.00 \\
TrQA Phi-3 & 0D$\to$LEN & 0.710 & 0.581 & -0.130 & [-0.156, -0.104] & no & 0.00 \\
TrQA Phi-3 & NONTOPO$\to$NONTOPO+DEFL & 0.883 & 0.885 & +0.002 & [+0.000, +0.004] & yes & 1.00 \\
TrQA Phi-3 & LEX$\to$LEX+0D & 0.635 & 0.721 & +0.086 & [+0.069, +0.104] & no & 0.00 \\
TrQA Phi-3 & LEX$\to$LEX+SINK & 0.635 & 0.718 & +0.083 & [+0.066, +0.101] & no & 0.00 \\
TrQA Phi-3 & LEX$\to$LEX+DEFL & 0.635 & 0.712 & +0.077 & [+0.057, +0.098] & no & 0.00 \\
TrQA Phi-3 & LOGPROB$\to$LOGPROB+0D & 0.877 & 0.878 & +0.000 & [-0.002, +0.002] & yes & 1.00 \\
TrQA Phi-3 & HIDDEN$\to$HIDDEN+0D & 0.867 & 0.867 & -0.000 & [-0.001, +0.000] & yes & 1.00 \\
TrQA Phi-3 & 0D$\to$SINK & 0.710 & 0.707 & -0.004 & [-0.010, +0.002] & yes & 0.98 \\
TrQA Phi-3 & 0D$\to$LOOKBACK & 0.710 & 0.840 & +0.129 & [+0.109, +0.150] & no & 0.00 \\
TrQA Phi-3 & NONTOPO$\to$NONTOPO+0D & 0.883 & 0.883 & -0.000 & [-0.001, +0.000] & yes & 1.00 \\
TrQA TinyLl. & LOGPROB$\to$LOGPROB+0D & 0.724 & 0.762 & +0.038 & [+0.025, +0.052] & no & 0.17 \\
TrQA TinyLl. & 0D$\to$SINK & 0.720 & 0.720 & +0.001 & [-0.010, +0.012] & yes & 0.40 \\
TrQA TinyLl. & LEX$\to$LEX+DEFL & 0.740 & 0.757 & +0.016 & [+0.002, +0.031] & no & 0.04 \\
TrQA TinyLl. & LEX$\to$LEX+SINK & 0.740 & 0.762 & +0.021 & [+0.010, +0.032] & no & 0.45 \\
TrQA TinyLl. & LEX$\to$LEX+0D & 0.740 & 0.764 & +0.024 & [+0.012, +0.036] & no & 0.33 \\
TrQA TinyLl. & NONTOPO$\to$NONTOPO+DEFL & 0.842 & 0.842 & +0.000 & [-0.003, +0.003] & yes & 1.00 \\
TrQA TinyLl. & NONTOPO$\to$NONTOPO+0D & 0.842 & 0.842 & -0.000 & [-0.001, +0.001] & yes & 1.00 \\
TrQA TinyLl. & 0D$\to$LEN & 0.720 & 0.631 & -0.089 & [-0.109, -0.068] & no & 0.00 \\
TrQA TinyLl. & 0D$\to$LEX & 0.720 & 0.740 & +0.021 & [-0.003, +0.044] & no & 0.00 \\
TrQA TinyLl. & 0D$\to$LLMCHECK & 0.720 & 0.688 & -0.032 & [-0.050, -0.013] & no & 0.00 \\
TrQA TinyLl. & 0D$\to$LOOKBACK & 0.720 & 0.823 & +0.104 & [+0.087, +0.121] & no & 0.00 \\
TrQA TinyLl. & HIDDEN$\to$HIDDEN+0D & 0.827 & 0.828 & +0.001 & [-0.001, +0.002] & yes & 1.00 \\
TrQA TinyLl. & 0D$\to$HIDDEN & 0.720 & 0.827 & +0.108 & [+0.088, +0.128] & no & 0.00 \\
TrQA TinyLl. & PH\_0D$\to$PH\_ENT & 0.775 & 0.807 & +0.032 & [+0.017, +0.047] & no & 0.01 \\
TrQA TinyLl. & PH\_SINK$\to$PH\_SINK+PH\_0DTOT & 0.814 & 0.813 & -0.001 & [-0.002, +0.000] & yes & 1.00 \\
TrQA TinyLl. & PH\_0DTOT$\to$PH\_SINK & 0.792 & 0.814 & +0.022 & [+0.015, +0.029] & no & 0.94 \\
TrQA TinyLl. & SINK$\to$SINK+ANS & 0.720 & 0.725 & +0.005 & [-0.011, +0.021] & no & 0.00 \\
TrQA TinyLl. & SINK$\to$SINK+1D & 0.720 & 0.719 & -0.001 & [-0.010, +0.008] & yes & 0.73 \\
TrQA TinyLl. & 0D$\to$0D+1D & 0.720 & 0.718 & -0.002 & [-0.011, +0.008] & yes & 0.60 \\
TrQA TinyLl. & SINK$\to$SINK+DEFL & 0.720 & 0.734 & +0.014 & [-0.000, +0.028] & no & 0.10 \\
TrQA TinyLl. & SINK$\to$0D+SINK & 0.720 & 0.724 & +0.004 & [-0.003, +0.011] & yes & 0.93 \\
TrQA TinyLl. & 0D$\to$LEGACY\_MSTPROXY & 0.720 & 0.729 & +0.010 & [-0.000, +0.020] & no & 0.55 \\
TrQA TinyLl. & 0D$\to$LOGPROB & 0.720 & 0.724 & +0.004 & [-0.019, +0.027] & no & 0.00 \\
TrQA TinyLl. & 0D$\to$ROWSTAT & 0.720 & 0.752 & +0.032 & [+0.020, +0.045] & no & 0.26 \\
\bottomrule
\end{longtable}
}
\normalsize

\section{Mistral-7B}
\label{app:mistral}
\MISTRALAPPENDIX

\section{Relation to accuracy-level audits, and cost}
\label{app:legacy}
Audits of attention-graph PH that stop at classification accuracy observe four facts that the theory now explains: (i) $0$D persistence equals an MST statistic (Fact~\ref{fact:mst}); (ii) $1$D persistence adds little beyond $0$D (Table~\ref{tab:tests}, column ``1D $\mid$ 0D''), which Theorem~\ref{thm:cone} now explains; (iii) the $1$D bank of Mistral-7B under its \texttt{[INST]} format is $99.99\%$ zeros, which is exactly the coned regime; (iv) HaluEval answers are separable from text alone.

\paragraph{Cost.} \TIMINGPARAGRAPH
